\chapter{Conclusiones}

\section{Conclusiones}

Este trabajo tenía como objetivo principal proponer una metodología de desarrollo de videojuegos que combinara los aspectos técnicos y creativos para posteriormente aplicarla en desarrollar un proyecto completo. \textbf{Se ha cumplido satisfactoriamente}: tras una revisión de la literatura disponible, se comprendieron los retos del desarrollo de videojuegos frente a otros tipos de software. A partir de esta comprensión se propuso una metodología de desarrollo que se empleó para producir un videojuego completo con éxito. El juego resultante puede jugarse desde el navegador \href{https://pavro-o.itch.io/mero-mero}{en este enlace} (contraseña: \textit{meromero}). A continuación se repasan los objetivos específicos.

\textbf{OE1}: \textit{Estudio de la literatura existente en los procesos de desarrollo en videojuegos. Entender sus retos y particularidades respecto a la ingeniería de software genérica}. Este objetivo se desarrolla exitosamente en el capítulo \ref{chap:stateofart}. Se comprendieron las dificultades del desarrollo desde el apartado técnico, artístico, de gestión y humano. Además, se indagó en otros temas como la gestión de calidad y las opciones de accesibilidad.

\textbf{OE2}: \textit{Realizar una planificación y gestión de un proyecto técnico abarcando las necesidades creativas de los videojuegos}. Se cumplió con éxito. En el capítulo \ref{chap:methodology} se expone una metodología que se implementa durante el capítulo \ref{chap:development} para producir el juego. Este se completó dentro del plazo de horas estimadas, por lo que se concluye que la planificación fue correcta. Adicionalmente, en el capítulo \ref{chap:methodology} se realiza una revisión de la metodología contrastándola con la experiencia obtenida durante el desarrollo.

\textbf{OE3}: \textit{Diseño e implementación de un juego de género Metroidvania usando el motor \textit{Godot}}. Se completó satisfactoriamente. Se produjo un juego completo de una duración aproximada de 30 minutos. En el capítulo \ref{chap:implementation} se desarrolla la implementación técnica del mismo.

\textbf{OE4}: \textit{Aplicación de mecanismos de gestión de calidad y validación para entregar un producto sólido, de calidad y que entretiene con éxito al usuario.} La naturaleza iterativa de la metodología usada, el enfoque a la planificación desde el prisma creativo, los mecanismos de gestión de calidad expuestos en el capítulo \ref{chap:implementation} y el extenso \textit{playtesting} que se documenta en el capítulo \ref{chap:development} (apartado de postproducción) ha permitido desarrollar un juego divertido y sólido que han podido jugar una veintena de personas. Aun así, se podría haber hecho más hincapié en el número de tests automáticos realizados, así como en la aplicación del \textit{snapshot testing} estudiado en el capítulo \ref{chap:stateofart}. Si no se hizo fue para poder dedicarle más tiempo a otras tareas durante la producción y cumplir con la planificación. Mencionar que aun con todo el juego no está libre de algunos \textit{bugs} conocidos, aunque no se interponen en la experiencia del jugador.

Para terminar, quisiera hacer una valoración personal sobre este trabajo. Para dar un poco de contexto, en mi tiempo libre programo y diseño videojuegos. Desde que empecé con esta afición, sentí mucha curiosidad por los procesos por los que se diseñan y producen videojuegos. Desde fuera, parecía casi un milagro que se llegaran a terminar algunos de estos títulos. Fue esta curiosidad la que me motivó a plantear este trabajo.

Hasta ahora había trabajado únicamente en proyectos pequeños, realizados en dos días o menos. Con este trabajo pude salir de mi zona de confort y afrontar un proyecto que abarcó varios meses. Estoy muy satisfecho por haber podido realizar un juego completo de estas dimensiones. He aprendido mucho sobre mi forma de trabajar y sobre cómo afrontar juegos de este calibre. Tras toda la investigación que realicé, me quedo con la sensación de que hay mucho camino que recorrer en la ingeniería aplicada a videojuegos, tanto a nivel personal como académico en general.

Irónicamente, no estoy de acuerdo con varias de las decisiones de diseño que yo mismo he tomado en el juego. Estas estaban influenciadas por el contexto de ser un trabajo académico y mi situación personal, compaginando trabajo y estudios. Pero, aunque el juego no sea perfecto, ver a mis amigos disfrutándolo y teorizando entre ellos para poder desentrañar sus secretos ha hecho que merezca la pena. Puedo asegurar que aplicaré lo aprendido en este trabajo en mis proyectos personales por venir.

\section{Trabajos futuros}

Se proponen las siguientes vías para trabajos futuros:

\begin{itemize}
    \item La continuación lógica de este trabajo sería realizar un segundo proyecto, aplicando la revisión de la metodología propuesta y evaluar su eficacia. Sería ideal contar con un perfil de artista en el equipo y así poder explorar la gestión de un equipo multidisciplinar.
    \item En esta memoria se comentó que el marco de trabajo Deep Design para prototipado de juegos funcionó bien con este proyecto, pero casaría mucho mejor en un juego más experimental en sus mecánicas, por lo que sería interesante realizar este segundo proyecto con un juego de ese perfil.
    \item Se propone aplicar un uso más extensivo de tests automáticos al proyecto, concretamente aplicando \textit{Test Driven Development} al desarrollo.
\end{itemize}
